%%=============================================================================
%% Inleiding
%%=============================================================================

\chapter{\IfLanguageName{dutch}{Inleiding}{Introduction}}%
\label{ch:inleiding}

%De inleiding moet de lezer net genoeg informatie verschaffen om het onderwerp te begrijpen en in te zien waarom de onderzoeksvraag de moeite waard is om te onderzoeken. In de inleiding ga je literatuurverwijzingen beperken, zodat de tekst vlot leesbaar blijft. Je kan de inleiding verder onderverdelen in secties als dit de tekst verduidelijkt. Zaken die aan bod kunnen komen in de inleiding~\autocite{Pollefliet2011}:

De toenemende digitalisering van bedrijfsprocessen en de groeiende afhankelijkheid van cloud- en netwerkinfrastructuren zorgen voor een exponentiële stijging in het aantal gegenereerde security-events. Security Operations Centers (SOC’s) worden dagelijks geconfronteerd met duizenden tot miljoenen meldingen afkomstig van logsystemen, firewalls, intrusion detection systemen en endpoint monitoring tools.
\newline
\newline
Hoewel Security Information and Event Management (SIEM)-systemen deze gebeurtenissen centraliseren en analyseren, blijft het een uitdaging om verbanden te leggen tussen deze grote hoeveelheden data. Moderne cyberaanvallen bestaan zelden uit één geïsoleerde actie, maar uit een keten van onderling gerelateerde gebeurtenissen. Het correct identificeren van deze verbanden is essentieel om complexe aanvalspatronen tijdig te detecteren.
\newline
\newline
In de praktijk worden veel security-events afzonderlijk beoordeeld of gecorreleerd via vooraf gedefinieerde regels. Deze rule-based aanpak werkt goed voor gekende aanvalspatronen, maar botst op beperkingen bij complexe, meerlagige aanvallen waarbij de relaties tussen events minder expliciet zijn. Dit leidt tot een hoge werkdruk voor SOC-analisten en verhoogt het risico op gemiste dreigingen. Bovendien zorgt het grote aantal alerts vaak voor zogenaamde “alert fatigue”, waarbij belangrijke signalen verloren kunnen gaan tussen minder relevante meldingen. Tegen deze achtergrond ontstaat de nood aan een meer contextuele en relationele benadering van dreigingsdetectie. Graph-based machine learning (GBML) biedt hier potentieel, omdat deze techniek relaties tussen entiteiten expliciet kan modelleren in de vorm van grafen. In plaats van events los te analyseren, worden ze voorgesteld als knooppunten (nodes) met onderlinge verbindingen (edges), waardoor aanvalsketens beter zichtbaar worden. 

%\begin{itemize}
 % \item context, achtergrond
  %\item afbakenen van het onderwerp
  %\item verantwoording van het onderwerp, methodologie
%  \item probleemstelling
 % \item onderzoeksdoelstelling
  %\item onderzoeksvraag
  %\item \ldots
%\end{itemize}

\section{\IfLanguageName{dutch}{Probleemstelling}{Problem Statement}}%
\label{sec:probleemstelling}

Voor cybersecuritybedrijven en hun SOC-analisten vormt het efficiënt correleren van security-events een concrete operationele uitdaging. Het manueel analyseren van grote hoeveelheden alerts is tijdrovend en foutgevoelig, terwijl klassieke SIEM-regels onvoldoende flexibel zijn om nieuwe of complexe aanvalspatronen correct te herkennen.
\newline
\newline
Het centrale probleem van deze bachelorproef is dan ook het gebrek aan een geautomatiseerde methode die security-events niet enkel individueel analyseert, maar ook hun onderlinge relaties kan leren en benutten om aanvalspatronen te reconstrueren.

%Uit je probleemstelling moet duidelijk zijn dat je onderzoek een meerwaarde heeft voor een concrete doelgroep. De doelgroep moet goed gedefinieerd en afgelijnd zijn. Doelgroepen als ``bedrijven,'' ``KMO's'', systeembeheerders, enz.~zijn nog te vaag. Als je een lijstje kan maken van de personen/organisaties die een meerwaarde zullen vinden in deze bachelorproef (dit is eigenlijk je steekproefkader), dan is dat een indicatie dat de doelgroep goed gedefinieerd is. Dit kan een enkel bedrijf zijn of zelfs één persoon (je co-promotor/opdrachtgever).

\section{\IfLanguageName{dutch}{Onderzoeksvraag}{Research question}}%
\label{sec:onderzoeksvraag}

Vanuit deze probleemstelling wordt de volgende hoofdonderzoeksvraag geformuleerd:
\newline
\newline
Hoe kan graph-based machine learning worden toegepast om security-events automatisch te correleren en te chainen, zodat complexe aanvalspatronen contextueler en overzichtelijker kunnen worden geanalyseerd binnen een SOC-context?
\newline
\newline
Deze onderzoeksvraag vertrekt niet vanuit het idee dat graph-based machine learning traditionele rule-based systemen volledig vervangt. De focus ligt eerder op het onderzoeken of een graph-based aanpak extra context kan toevoegen aan bestaande detectie- en correlatiemechanismen.
\newline
\newline
Om deze vraag te beantwoorden, wordt een onderscheid gemaakt tussen het probleemdomein en het oplossingsdomein.
\newline
\newline
\textbf{Binnen het probleemdomein wordt nagegaan:}
\begin{itemize}
    \item Welke types security-events en datavelden relevant zijn om verbanden tussen events te leggen.
    \item Waarom klassieke rule-based correlatie moeilijk schaalbaar is bij grote hoeveelheden alerts.
    \item Welke rol tijd, assets, IP-adressen, gebruikers en eventtypes spelen bij het reconstrueren van mogelijke aanvalsketens.
    \item Waarom labels, class imbalance en heterogene dataformaten een uitdaging vormen binnen cybersecuritydatasets.
\end{itemize}

\textbf{Binnen het oplossingsdomein wordt onderzocht:}
\begin{itemize}
    \item Hoe security-events kunnen worden omgezet naar een grafenstructuur met nodes en edges.
    \item Hoe een Relational Graph Convolutional Network gebruikt kan worden om events te classificeren of te scoren.
    \item Hoe modeloutput gebruikt kan worden om verdachte events te rangschikken.
    \item Hoe verdachte events via een chaining-stap gegroepeerd kunnen worden op basis van gedeelde entiteiten en temporele nabijheid.
    \item Welke beperkingen optreden wanneer labels onvolledig zijn of wanneer de dataset sterk onevenwichtig verdeeld is.
\end{itemize}
%Wees zo concreet mogelijk bij het formuleren van je onderzoeksvraag. Een onderzoeksvraag is trouwens iets waar nog niemand op dit moment een antwoord heeft (voor zover je kan nagaan). Het opzoeken van bestaande informatie (bv. ``welke tools bestaan er voor deze toepassing?'') is dus geen onderzoeksvraag. Je kan de onderzoeksvraag verder specifiëren in deelvragen. Bv.~als je onderzoek gaat over performantiemetingen, dan 

\section{\IfLanguageName{dutch}{Onderzoeksdoelstelling}{Research objective}}%
\label{sec:onderzoeksdoelstelling}

Het doel van deze bachelorproef is het ontwikkelen en evalueren van een proof-of-concept waarin security-events worden voorgesteld als een graafstructuur en geanalyseerd met graph-based machine learning.
\newline
\newline
Binnen deze proof-of-concept worden twee experimentele settings uitgewerkt. Eerst wordt een supervised multiclass classificatie-experiment uitgevoerd op CIC-IDS2018. Deze dataset bevat duidelijke labels en laat toe om na te gaan of een R-GCN-model verschillende aanvalstypes kan onderscheiden.
\newline
\newline
Daarna wordt dezelfde graph-based aanpak toegepast op LANL-data. Deze dataset bevat meer realistische enterprise-events, maar de labels zijn beperkter en worden afgeleid uit redteam-informatie. Hierdoor wordt deze setting niet behandeld als een zuivere multiclass classificatietaak, maar als een weakly supervised suspicious-event detection probleem waarbij ranking en chaining centraal staan.
\newline
\newline
Het prototype moet aantonen in welke mate deze aanpak:
\begin{itemize}
    \item security-events kan voorstellen als een relationele graaf;
    \item aanvalstypes kan onderscheiden wanneer duidelijke labels beschikbaar zijn;
    \item verdachte events kan rangschikken wanneer labels beperkter of zwakker zijn;
    \item events kan groeperen tot mogelijke chains op basis van tijd en gedeelde entiteiten;
    \item inzicht geeft in de beperkingen van graph-based machine learning bij class imbalance en weak labels.
\end{itemize}

De performantie wordt geëvalueerd met metrics die aansluiten bij de rol van elke dataset. Voor CIC-IDS2018 worden klassieke classificatiemetrics gebruikt, zoals accuracy, precision, recall en F1-score. Voor LANL wordt daarnaast gekeken naar ranking- en chaininggerichte evaluatie, omdat deze beter aansluit bij de manier waarop verdachte events in een SOC-context onderzocht worden.

\section{\IfLanguageName{dutch}{Opzet van deze bachelorproef}{Structure of this bachelor thesis}}%
\label{sec:opzet-bachelorproef}

% Het is gebruikelijk aan het einde van de inleiding een overzicht te
% geven van de opbouw van de rest van de tekst. Deze sectie bevat al een aanzet
% die je kan aanvullen/aanpassen in functie van je eigen tekst.
De rest van deze bachelorproef is als volgt opgebouwd:

In Hoofdstuk~\ref{ch:stand-van-zaken} wordt een overzicht gegeven van de stand van zaken binnen het onderzoeksdomein, op basis van een literatuurstudie.

In Hoofdstuk~\ref{ch:methodologie} wordt de methodologie toegelicht en worden de gebruikte onderzoekstechnieken besproken om een antwoord te kunnen formuleren op de onderzoeksvragen.

% TODO: Vul hier aan voor je eigen hoofstukken, één of twee zinnen per hoofdstuk

In Hoofdstuk~\ref{ch:conclusie}, tenslotte, wordt de conclusie gegeven en een antwoord geformuleerd op de onderzoeksvragen. Daarbij wordt ook een aanzet gegeven voor toekomstig onderzoek binnen dit domein.